% Originally: the \section{Solutions} of content/week-05.tex and content/week-06.tex

\section{Solutions}
\label{sec:extrema_solutions}

\begin{exercisesolution}[Proof of \cref{ex:ai_extrema_triangle}]
% Generator: Gemini 3.6 Flash (Medium)
Since $T$ is compact and $f$ is continuous, global extrema exist by the extreme value theorem
(\cref{item:extreme_value_theorem}).
\begin{itemize}
  \item \textbf{Interior of $T$:} Critical points satisfy $\nabla f(x,y) = (y,x)\transp = (0,0)\transp$, and therefore $(x,y) = (0,0)$, which lies on the boundary. Thus, there are no critical points in $\operatorname{int}(T)$.
  \item \textbf{Boundary of $T$:}
    \begin{itemize}
      \item Along the leg $x = 0$ ($0 \le y \le 2$), $f(0,y) = 0$.
      \item Along the leg $y = 0$ ($0 \le x \le 2$), $f(x,0) = 0$.
      \item Along the hypotenuse $x+y = 2$ ($0 \le x \le 2$), substitute $y = 2-x$:
        \[h(x) := x(2-x) = 2x-x^2, \qquad h'(x) = 2-2x,\]
        which vanishes at $x = 1$. This gives the candidate point $(1,1)$ with $f(1,1) = 1$. The endpoints $(2,0)$ and $(0,2)$ yield $f = 0$.
    \end{itemize}
\end{itemize}
Comparing values, the \textbf{global maximum} is $1$ at $(1,1)$, and the \textbf{global minimum} is $0$, attained everywhere on the coordinate axes legs of $T$.
\end{exercisesolution}

\begin{exercisesolution}[Proof of \cref{ex:6.2}]
% Generator: Claude Sonnet 5 (Medium)
By the spectral theorem, since $M$ is symmetric there exists an orthogonal matrix $O$
(so $OO\transp = \id$, hence $\det O = \pm 1$) and $\lambda_1, \lambda_2 \in \mathbb{R}$ such that
\[OMO\transp = \begin{pmatrix} \lambda_1 & 0 \\ 0 & \lambda_2 \end{pmatrix}.\]
Using $\det(AB) = \det(A)\det(B)$ and $\Tr(AB) = \Tr(BA)$,
\[\det M = \det(OMO\transp) = \lambda_1\lambda_2, \qquad \Tr M = \Tr(MO\transp O) = \Tr(OMO\transp) = \lambda_1 + \lambda_2.\]
So, $\det M$ and $\Tr M$ are exactly the product and sum of the eigenvalues, and the four cases
follow immediately:
\begin{itemize}
  \item $\det M > 0$: $\lambda_1, \lambda_2$ have the same sign. If additionally $\Tr M > 0$,
    both are positive, so $M$ is \textbf{positive definite}; if $\Tr M < 0$, both are negative,
    so $M$ is \textbf{negative definite}.
  \item $\det M < 0$: $\lambda_1, \lambda_2$ have opposite signs, so $M$ is \textbf{indefinite}.
  \item $\det M = 0$: one of $\lambda_1, \lambda_2$ is zero, so $M$ is \textbf{degenerate}.
\end{itemize}
\end{exercisesolution}

\begin{exercisesolution}[Solution to \cref{ex:6.8}]
% Generator: Claude Sonnet 5 (Medium)
None of \textbf{(a)}--\textbf{(c)} hold in general:
\begin{itemize}
  \item \textbf{(a) False:} $f \equiv 0$ has $\Hess f(x_0) = 0$ positive semidefinite at every
    $x_0$, but $x_0$ is not a \emph{strict} local minimum (every nearby value equals $f(x_0)$).
  \item \textbf{(b) False:} $f(x) = x_1^3$ at $x_0 = 0$ has $\Hess f(0) = 0$ (positive
    semidefinite), but $f$ takes negative values arbitrarily close to $0$, so $x_0$ is not a
    local minimum at all.
  \item \textbf{(c) False:} $f(x) = -x_1^4$ at $x_0 = 0$ has $\Hess f(0) = 0$ (positive
    semidefinite, vacuously, since it is the zero matrix), yet $x_0$ \emph{is} a (strict) local
    maximum of $f$.
\end{itemize}
Since \textbf{(a)}--\textbf{(c)} all fail, \textbf{(d) is the correct answer}: none of the above
statements necessarily hold.
\end{exercisesolution}

\begin{exercisesolution}[Solution to \cref{ex:6.9}]
% Generator: Claude Sonnet 5 (Medium)
\begin{enumerate}[label=\textbf{(\alph*)}]
  \item With $g(x,y) := x^2+y^2-1$, the Lagrangian is $L(x,y,\lambda) := f(x,y) - \lambda g(x,y)$,
    giving the system
    \[0 = \partial_x L = 2x(2-\lambda) - 1, \qquad 0 = \partial_y L = 2y(2-2\lambda), \qquad 0 = \partial_\lambda L = -(x^2+y^2-1).\]
    From the middle equation, either $y = 0$ or $\lambda = 1$.

    \textbf{If $y = 0$:} then $x^2 = 1$, so $x = \pm 1$, giving the points $(1,0)$ and $(-1,0)$
    with $f(1,0) = 1$ and $f(-1,0) = 3$.

    \textbf{If $\lambda = 1$:} the first equation gives $2x(1) = 1$, i.e.\ $x = \tfrac12$, and the
    constraint gives $y = \pm\tfrac{\sqrt3}{2}$, with $f\big(\tfrac12,\pm\tfrac{\sqrt3}{2}\big) =
    \tfrac12 + \tfrac34 - \tfrac12 = \tfrac34$.

    Comparing all four values, $f$ attains a \textbf{global maximum} at $(-1,0)$ (value $3$) and
    \textbf{global minima} at $\big(\tfrac12,\pm\tfrac{\sqrt3}{2}\big)$ (value $\tfrac34$); $(1,0)$
    is a local (non-global) extremum with value $1$.
  \end{enumerate}
\end{exercisesolution}

\begin{ainote}
The official solution lists $(1,0)$ alongside $(-1,0)$ as \qt{local extrema} from the Lagrange
system, then only mentions it again implicitly by omission from the final global-extrema summary.
It is worth flagging explicitly that \textbf{Lagrange's condition finds all \emph{critical} points
of $f|_{S^1}$, not automatically the global extrema}. Here $(1,0)$ is a genuine local maximum along
the circle (value $1$) that is simply not the largest one; only comparing all candidate values
at the end distinguishes it from the global maximum at $(-1,0)$ (value $3$).
\end{ainote}

\begin{exercisesolution}[Solution to \cref{ex:6.9}, continued]
% Generator: Claude Sonnet 5 (Medium)
\begin{enumerate}[label=\textbf{(\alph*)}, start=2]
  \item Having handled the boundary $S^1 = \partial D$ in \textbf{(a)}, we look for interior
    critical points: $Df(x,y) = (4x-1, 2y) = 0$ gives $(x,y) = \big(\tfrac14, 0\big) \in
    \operatorname{int}(D)$. The Hessian $\Hess f = \left(\begin{smallmatrix} 4 & 0 \\ 0 & 2\end{smallmatrix}\right)$ is
    positive definite, so this is a local minimum, with value $f\big(\tfrac14,0\big) =
    \tfrac18 - \tfrac14 = -\tfrac18$. Since $-\tfrac18$ is smaller than every boundary value found
    in \textbf{(a)}, $f$ on $D$ attains a \textbf{global minimum} at $\big(\tfrac14,0\big)$ and a
    \textbf{global maximum} at $(-1,0)$.
\end{enumerate}
\end{exercisesolution}

\begin{exercisesolution}[Solution to \cref{ex:6.12}]
% Generator: Claude Sonnet 5 (Medium)
Writing $r^2 := x^2+y^2$, the gradient is
\[\nabla f(x,y) = \big(2x(a-ax^2-by^2),\ 2y(b-ax^2-by^2)\big)e^{-r^2}.\]

\textbf{Case $a=b=0$:} $f \equiv 0$, so every point of $\mathbb{R}^2$ is critical, and (being the
constant function) each is simultaneously a global maximum and minimum.

\textbf{Case $a=0, b\neq 0$ (and symmetrically $a \neq 0, b=0$):} the gradient equations reduce to
$-2xby^2=0$ and $2by(1-y^2)=0$, whose common solutions are the entire line $\{(x,0): x\in
\mathbb{R}\}$ together with the two points $(0,\pm1)$. On the critical line the Hessian is
singular (degenerate, as expected since a whole line is critical); at $(0,\pm1)$,
$\Hess f(0,\pm1) = \operatorname{diag}(-2b,-4b)\,e^{-1}$, giving \textbf{minima} if $b<0$ and
\textbf{maxima} if $b>0$. The case $a\neq0,b=0$ is symmetric, with critical line $\{(0,y)\}$ and
extra points $(\pm1,0)$: minima if $a<0$, maxima if $a>0$.

\textbf{Case $a,b \neq 0$, $a \neq b$:} setting $x=0$ or $y=0$ in the gradient equations gives the
five critical points $(0,0)$, $(0,\pm1)$, $(\pm1,0)$ (checking $x,y\neq0$ simultaneously forces
$a=b$, excluded here). The Hessians are
\begin{align*}
\Hess f(0,0) &= \begin{pmatrix}2a&0\\0&2b\end{pmatrix}, \\
\Hess f(0,\pm1) &= e^{-1}\begin{pmatrix}2(a-b)&0\\0&-4b\end{pmatrix}, \\
\Hess f(\pm1,0) &= e^{-1}\begin{pmatrix}-4a&0\\0&2(b-a)\end{pmatrix}.
\end{align*}
Reading off signs for each ordering of $a,b$ against $0$:
\begin{itemize}
  \item $a>b>0$: min at $(0,0)$, saddle at $(0,\pm1)$, max at $(\pm1,0)$.
  \item $b>a>0$: min at $(0,0)$, max at $(0,\pm1)$, saddle at $(\pm1,0)$.
  \item $a>0>b$: saddle at $(0,0)$, min at $(0,\pm1)$, max at $(\pm1,0)$.
  \item $b>0>a$: saddle at $(0,0)$, max at $(0,\pm1)$, min at $(\pm1,0)$.
  \item $0>a>b$: max at $(0,0)$, min at $(0,\pm1)$, saddle at $(\pm1,0)$.
  \item $0>b>a$: max at $(0,0)$, saddle at $(0,\pm1)$, min at $(\pm1,0)$.
\end{itemize}

\textbf{Case $a=b\neq0$:} beyond $(0,0)$, $(0,\pm1)$, $(\pm1,0)$, setting $x,y\neq0$ now gives
$a-ax^2-ay^2=0$, that is, $x^2+y^2=1$, so the entire unit circle is critical. On that circle,
\[\Hess f(x,y) = -4a\,e^{-1}\begin{pmatrix} x^2 & xy \\ xy & y^2 \end{pmatrix},\]
whose determinant is $16a^2e^{-2}(x^2y^2 - x^2y^2) = 0$. The Hessian is therefore degenerate
(\cref{def:sign_symmetric_matrix}) at every such point, and the test is inconclusive, as it must
be, since a whole curve of critical points cannot consist of isolated extrema.
\end{exercisesolution}

\begin{exercisesolution}[Solution to \cref{ex:ai_fta_hypotheses}]
% Generator: Claude Opus 5 (High)
\begin{enumerate}[label=\textbf{(\alph*)}]
  \item The restriction is what makes $\lvert z\rvert^k \leq \lvert z\rvert^{n-1}$ for every
    $k \leq n-1$, which is the step that collapses the whole sum into a single power. For
    $\lvert z\rvert < 1$ the inequality reverses, the low powers dominate, and the estimate is
    simply false: at $z = 0$ it would read $\lvert a_0\rvert \geq 0 \cdot (0 - C)$ only by
    accident of the right-hand side being negative, and it gives nothing usable.

    Nothing is lost, because the estimate is only ever applied for $\lvert z\rvert \geq R$, and
    $R \geq 1$ automatically: $C = \sum_{k=0}^n\lvert a_k\rvert$ includes the term
    $\lvert a_n\rvert = 1$, so $C \geq 1$ and hence $R \geq C+1 \geq 2$. The hypothesis is free.
  \item Here $f(z) = \lvert z\rvert^2+1 \geq 1$, so $f$ has no root, and indeed
    $\lvert f\rvert$ attains its minimum $1$ at $z_0 = 0$. Both halves of the first part of the
    proof therefore go through unchanged: $\lvert f\rvert$ grows without bound, the infimum is
    attained, and the minimiser is interior.

    The step that breaks is the expansion. The proof writes $g(z) := f(z_0+z) = \sum_k b_kz^k$,
    and $f$ here is not a polynomial in $z$: it is $z\overline z + 1$, and $\overline z$ is not
    expressible in powers of $z$. Concretely, at $z_0 = 0$ and $z = re^{i\varphi}$,
    \[g\big(re^{i\varphi}\big) = 1 + r^2,\]
    in which the perturbation carries \emph{no} factor $e^{i\ell\varphi}$. The direction
    $\varphi$ has vanished from the expression, so there is no direction to aim, and the term
    $r^2$ increases $\lvert g\rvert$ whichever way one leaves the origin.

    That is the sharpest way to say what the proof actually uses: not that $f$ is smooth, nor
    that it grows, but that the lowest-order correction to $b_0$ has a \emph{phase} that the free
    parameter $\varphi$ can rotate onto $-b_0$. The function $z\overline z$ is the standard
    example of a perturbation with no phase to rotate.
  \item The requirement is $e^{i(\ell\varphi+\psi)} = -1$, i.e.\
    $\ell\varphi + \psi \equiv \pi \pmod{2\pi}$, so
    \[\varphi_j := \frac{\pi - \psi + 2\pi j}{\ell}, \qquad j = 0,1,\ldots,\ell-1,\]
    which are $\ell$ distinct directions in $[0,2\pi)$, equally spaced by $2\pi/\ell$. The
    argument does not care which is used, because it needs only \emph{one} ray along which
    $\lvert g\rvert$ falls below $\lvert b_0\rvert$; the contradiction is reached from a single
    such ray. That $\ell$ of them exist rather than one is a bonus of working over $\mathbb{C}$,
    and it is exactly the freedom that \textbf{(b)} shows to be indispensable.
\end{enumerate}
\end{exercisesolution}

\begin{exercisesolution}[Proof of \cref{ex:hessian_parameterized_exp}]
% Generator: Gemini 3.6 Flash (High)
% Correction: Claude Opus 5 (High) --- Hf -> \Hess f, \text{Tr} -> \Tr, ^T -> \transp, and the
% \implies signs standing in for verbs replaced by prose. The range argument in (c) also spelled
% out: "by continuity" was doing a lot of unstated work, and it is the connectedness of R^2 plus
% the intermediate value theorem that closes the gaps.
\begin{enumerate}[label=\textbf{(\alph*)}]
  \item The gradient of $f_\alpha$ is
\[\nabla f_\alpha(x,y) = \begin{pmatrix} e^x(x^2 + \alpha y^2 + 2x) \\ 2\alpha y e^x \end{pmatrix},\]
and we look for the points where it vanishes. Since $e^x > 0$ everywhere, the second component
vanishes when $\alpha y = 0$, which splits the problem in two.
\begin{itemize}
  \item \textbf{If $\alpha = 0$:} the second component vanishes identically, and the first reduces
    to $x^2 + 2x = 0$, so $x = 0$ or $x = -2$. The critical set is the pair of vertical lines
    $\{(0,y) \mid y \in \mathbb{R}\}$ and $\{(-2,y) \mid y \in \mathbb{R}\}$.
  \item \textbf{If $\alpha \neq 0$:} the second component forces $y = 0$, and the first then reads
    $x^2 + 2x = 0$ again. The critical points are $(0,0)$ and $(-2,0)$.
\end{itemize}

  \item For $\alpha \neq 0$, differentiating once more gives
\[\Hess f_\alpha(x,y) = e^x \begin{pmatrix} x^2+4x+2+\alpha y^2 & 2\alpha y \\ 2\alpha y & 2\alpha \end{pmatrix}.\]
\begin{itemize}
  \item \textbf{At $(0,0)$:} $\Hess f_\alpha(0,0) = \left(\begin{smallmatrix} 2 & 0 \\ 0 & 2\alpha \end{smallmatrix}\right)$, with determinant $4\alpha$. For $\alpha > 0$ the determinant is positive and $\Tr = 2+2\alpha > 0$, so the matrix is positive definite and $(0,0)$ is a strict local minimum. For $\alpha < 0$ the determinant is negative, so $(0,0)$ is a saddle point.
  \item \textbf{At $(-2,0)$:} $\Hess f_\alpha(-2,0) = e^{-2}\left(\begin{smallmatrix} -2 & 0 \\ 0 & 2\alpha \end{smallmatrix}\right)$, with determinant $-4\alpha e^{-4}$. For $\alpha > 0$ this is negative, so $(-2,0)$ is a saddle point. For $\alpha < 0$ it is positive, and the entry $\partial_1^2 f_\alpha(-2,0) = -2e^{-2} < 0$ makes the matrix negative definite, so $(-2,0)$ is a strict local maximum.
\end{itemize}
Notice that the two critical points swap roles as $\alpha$ changes sign, and that neither test is
inconclusive: $\det \Hess$ vanishes only at $\alpha = 0$, the case excluded here.

  \item Each range is determined by exhibiting the extreme values and then filling in between by the
    intermediate value theorem, which applies because $\mathbb{R}^2$ is connected and $f_\alpha$ is
    continuous.
\begin{itemize}
  \item \textbf{If $\alpha > 0$:} both $x^2$ and $\alpha y^2$ are non-negative, so $f_\alpha \geq 0$,
    with $f_\alpha(0,0) = 0$ attained. Along the $x$-axis, $f_\alpha(x,0) = x^2 e^x \to \infty$ as
    $x \to \infty$, so values grow without bound. Hence $f_\alpha(\mathbb{R}^2) = [0,\infty)$.
  \item \textbf{If $\alpha < 0$:} along the $x$-axis $f_\alpha(x,0) = x^2 e^x$ sweeps out
    $[0,\infty)$ as above, while along the $y$-axis $f_\alpha(0,y) = \alpha y^2$ sweeps out
    $(-\infty,0]$. Together these give $f_\alpha(\mathbb{R}^2) = \mathbb{R}$.
  \item \textbf{If $\alpha = 0$:} $f_0(x,y) = x^2 e^x$ does not depend on $y$ and is non-negative,
    with value $0$ at $x = 0$ and unbounded above. Hence $f_0(\mathbb{R}^2) = [0,\infty)$.
\end{itemize}
\end{enumerate}
\begin{ainote}
It is worth pausing on the $\alpha > 0$ case: $(0,0)$ is a strict local minimum \emph{and} the global
minimum, but $(-2,0)$ is only a saddle, and $f_\alpha$ has no global maximum at all. Part
\textbf{(c)} is not a corollary of part \textbf{(b)} --- the Hessian test is purely local, and
reading a range off it is a standard way to go wrong.
\end{ainote}
\end{exercisesolution}


